\section{Introduction}

The previous chapter reported the empirical results of PatchOps. This chapter discusses what those results mean
for the design of AI agents that operate inside CI pipelines. The central point is that PatchOps should not be
read as an attempt to maximize patch volume at any cost. It is better understood as a system that makes repair
conditional: act when the failure signal, repository evidence, and safety checks support action; otherwise
abstain, preview, rerun, or escalate.

This distinction is important because CI repair is an unusually tempting domain for overclaiming. A failed CI
run provides an apparently crisp objective, and a green rerun appears to offer a simple success signal. The
experiments in this thesis show that the real problem is less tidy. The failure log may identify a symptom
rather than the file that should be changed. A patch may apply cleanly but fail semantically. A green CI result
may still be produced by an off-target or dishonest edit. An agent may be given the option to abstain and still
rarely choose it. These observations are why the safety-gated and abstention-aware framing is not a decorative
addition to PatchOps; it is the core design response to the empirical behavior of the task.

\section{Why the Negative Results Matter}

Several of the most important results in PatchOps are negative results. Full-repository lexical retrieval did
not break the localization ceiling. LLM reranking did not fix whole-repository retrieval. Semantic embeddings
did not outperform the fused lexical baseline. Bounded agentic exploration did not reliably beat cheaper
localizers. Full-repository retrieval made generated patches apply, but did not make the gated system achieve
CI Pass@1 on the measured repair subset.

These results could be read superficially as a sequence of failed improvements. In a thesis, however, their
value is different. They narrow the problem. Before these experiments, low repair success could have been
explained by weak context, missing files, invalid diffs, poor classification, or lack of repository search. The
later phases remove several of those explanations. Classification becomes competitive with an LLM through a
local learned model. Full-repository retrieval removes the apply-rate bottleneck among attempted gated patches.
Embeddings and agentic navigation test the obvious next localization levers. The remaining wall is therefore
more specific: semantic repository reasoning and repair correctness.

This sharpening is a real research contribution. A system that only reports end-to-end repair success can hide
where the failure occurs. PatchOps decomposes the pipeline into diagnosis, localization, patch production,
patch application, validation, safety, and abstention. That decomposition makes it possible to say not only that
CI repair is hard, but where and why it remains hard in this implementation.

\section{PatchOps as a Conditional Repair System}

The most important design implication is that CI repair agents should be conditional systems rather than
patch-emission systems. A patch-emission system treats every failure as an opportunity to generate a diff. A
conditional repair system treats patching as only one possible decision among several. PatchOps therefore
distinguishes between \texttt{patch}, \texttt{preview}, \texttt{rerun}, and \texttt{escalate}.

This decision vocabulary changes the meaning of success. A correct rerun decision for a flaky failure is not a
failure to repair; it is the right action. Escalating a secret, permission, deployment, or ambiguous failure is
not a weakness if automatic editing would be unsafe. Previewing a low-plausibility patch is not indecision if
the available evidence does not justify opening an automatic PR. In this framing, abstention is not the absence
of intelligence. It is one of the system's safety behaviors.

The results support this design. The labeled safety set shows that explicit decision classes can be evaluated.
The selective-prediction experiment shows that abstention can be converted into a risk-coverage trade-off
instead of a vague rule. The patch-plausibility experiment shows that even after a patch exists, the system can
separate automatic action from preview. The reward-hacking detector shows that a patch can be rejected because
it is dishonest, not merely because it fails to apply.

\section{Why Abstention Must Be Engineered}

One of the clearest lessons from the agentic localization study is that abstention cannot simply be delegated to
the base model. The bounded repository-exploration agent had an explicit abstain action, but it used it rarely
and run-dependently. When it did abstain, the abstentions were justified, but the frequency was too unstable to
serve as the main safety mechanism.

This result explains the need for PatchOps' external decision stack. The learned classifier confidence provides
a stronger diagnosis-level abstention signal than the hand-written engine confidence. Phase 15 then turns that
calibrated signal into a learned act-versus-hold gate for the low-risk lane, replacing the older hand-set
threshold decision. Reward-hacking detection adds a hard veto for dishonest patches. Patch plausibility provides
a patch-level signal for whether a diff should be auto-PR'd or merely previewed. Together, these components make
abstention a property of the system architecture rather than a hopeful behavior of the LLM.

This distinction also matters for accountability. If an LLM freely decides when to decline, the system has only
a weak explanation for why it acted on one failure and not another. If abstention is tied to calibrated
confidence, a learned act-versus-hold model, safety predicates, and patch-locality signals, the system can
produce a traceable reason: low diagnosis confidence, high-risk edit target, reward-hacking pattern, or low
plausibility. That traceability is essential for a CI agent that may open PRs in a real repository.

\section{CI-Green as a Partial Oracle}

CI validation is valuable, but it is not a complete correctness oracle. PatchOps uses CI re-execution because it
is the strongest readily available feedback signal in this domain. If a generated fix still leaves CI red, the
system has direct evidence that the repair failed. However, the reverse does not always hold. A green run may
mean the underlying problem was fixed, but it may also mean the check was weakened, the failing test was removed,
or the patch happened to satisfy the current test suite while missing the semantic intent.

This is why the thesis connects CI repair to the broader APR overfitting literature. The risk is not only that
an agent fails to fix a bug. The more subtle risk is that it produces a plausible patch that optimizes the
observable validation signal without respecting the intended behavior. PatchOps' reward-hack and plausibility
gates are responses to that risk. They do not prove semantic correctness, but they make the system less willing
to treat every green or syntactically valid patch as acceptable.

The phrase ``CI-validated repair'' should therefore be read carefully. It means that selected fixes are checked
against CI, not that CI fully certifies correctness. The thesis contribution is the combination of CI validation
with safety and abstention controls, not a claim that CI alone is enough.

\section{The Role of the Learned Classifier}

The learned classifier has a deliberately modest role in the thesis. It is not presented as the main technical
breakthrough and it is not claimed to improve CI Pass@1. Its value is practical and methodological. It makes one
stage of the pipeline cheaper, reproducible, and trainable from public data while matching the LLM's observed
classification band on the measured split.

This matters because CI agents are often evaluated under conditions that are difficult to reproduce exactly:
closed models, drifting APIs, hidden prompts, and paid inference. Moving classification onto a small local model
does not remove all of those issues, because the rest of PatchOps still uses LLMs for generation and reasoning.
However, it reduces reliance on closed inference for a repeated pipeline stage and provides a stable confidence
signal for selective prediction.

The classifier is therefore best understood as infrastructure for trustworthy operation. It supports the
abstention story, improves cost behavior, and strengthens reproducibility. It does not distract from the main
finding that localization and semantic repair remain the hard parts.

\section{The Learned Decision Gate as the Abstention Capstone}

Phase 15 resolves a weakness that an examiner would naturally notice in the earlier architecture: if the system
uses confidence thresholds, where did those thresholds come from? The answer in the final PatchOps artifact is
that the low-risk act-versus-hold decision is no longer only a hand-set threshold. It is an optional learned,
calibrated decision gate trained on the labeled Safety Set.

This matters because it converts the phrase ``abstention must be engineered'' into an implemented mechanism. The
gate uses signals available before patch generation and produces a calibrated probability of action. On the
expanded 76-case Safety Set it matches-or-beats the hand-tuned threshold policy, with overlapping confidence
intervals, and on the low-risk lane it recovers confident patch actions that the old thresholds held back. The
correct discussion point is not that PatchOps suddenly repairs more CI failures end to end. The point is that the
act-versus-hold boundary is now measured, calibrated, and inspectable.

The result also preserves the conservative structure of PatchOps. The learned gate does not bypass flaky
reruns, high-risk escalation, reward-hack detection, or patch plausibility. It only decides whether a low-risk
case should proceed toward action or remain held for preview. This is the right level of learning for the thesis
claim: use learned calibration to remove brittle constants, but keep deterministic safety boundaries in place.

\section{Architecture: Why a Staged Pipeline Is Preferable Here}

PatchOps uses a staged decision pipeline rather than an unconstrained multi-agent architecture. The experiments
support that choice. In CI repair, each stage has a different failure mode: diagnosis can misclassify the
failure, localization can retrieve the wrong file, patch generation can overfit, CI validation can be gamed, and
the planner can act too aggressively. A staged architecture lets the system attach checks and evidence to each
boundary.

The agentic localization phase is especially informative here. Tool use is attractive because it resembles how
developers investigate repositories: search, read, follow imports, and decide. In the measured setting, however,
the bounded agent was slower than deterministic retrieval and did not provide a global localization gain. That
does not mean agentic exploration is useless; the clean wins show that it can help in some cases. It does mean
that agentic behavior should be introduced as a bounded component with fallback and external gates, not as an
unrestricted replacement for the whole pipeline.

The practical design lesson is that AI CI repair benefits from explicit control surfaces. Deterministic
retrieval, learned classification, the learned decision gate, LLM reasoning, safety gates, plausibility checks,
and validation runners each contribute different kinds of evidence. The thesis argues for orchestrating these
pieces conservatively rather than trusting a single model call to diagnose, localize, patch, validate, and
self-regulate.

\section{Implications for Deployment}

In a real engineering team, PatchOps would be most appropriate as a tiered assistant rather than a fully
autonomous maintainer. For low-risk, log-localizable failures such as straightforward version bumps or narrowly
scoped dependency issues, the system can propose or open a small PR after safety checks. For ambiguous failures,
test-semantics failures, permissions issues, deployment failures, and low-plausibility patches, the system
should provide a preview or escalation with evidence instead of changing code automatically.

This operating model is less glamorous than a universal auto-fixer, but it is more realistic. CI pipelines are
part of the trust boundary of a repository. A bad patch can hide a real regression, weaken tests, or modify
workflow behavior in ways that have security and release consequences. The correct default is therefore not
``always patch.'' The correct default is ``patch only when the evidence and gates justify it.''

The results also suggest how teams could tune PatchOps. A team that prioritizes speed could lower abstention
thresholds or the learned gate's operating point and accept more previewed patches. A team that manages critical
infrastructure could require higher plausibility and confidence before auto-PR. The important point is that
PatchOps exposes this as a measurable coverage-versus-risk trade-off rather than as a hidden prompt preference.

\section{Relation to Prior CI Repair Work}

The discussion also clarifies how PatchOps should be positioned relative to CI-Repair-Bench and related repair
systems. PatchOps does not claim to be the first system to perform CI-verified repair, and it does not claim to
beat the strongest reported repair baselines on raw Pass@1. Those would be the wrong claims. CI-Repair-Bench
already establishes repository-level CI-verified repair as a benchmarked task, and its strongest baseline sets a
clear reference point.

PatchOps contributes along a different axis. It asks what an agent should do when repair evidence is weak, when
patches are plausible but off-target, when the validation oracle can be gamed, and when the model does not
reliably abstain on its own. This safety-and-abstention axis is not a substitute for repair accuracy; it is a
complement to it. In practical CI automation, both are needed. A high-repair system that cannot distinguish safe
from unsafe action is risky, while a safe system with low coverage is useful only on a narrower subset. PatchOps
occupies the second side of that trade-off and makes the trade-off explicit.

\section{Chapter Conclusion}

This chapter interpreted the empirical results of PatchOps as evidence for conditional, safety-gated CI repair.
The main discussion points are that negative results help isolate the true bottleneck, abstention must be
engineered rather than entrusted to the model alone, Phase 15 turns the key low-risk act-versus-hold boundary
into a learned calibrated gate, CI-green is only a partial correctness signal, and the staged architecture is
justified by the distinct failure modes of diagnosis, localization, generation, validation, and decision-making.

The next chapter states the limitations and threats to validity of these claims. It is intentionally explicit
about scope, sample sizes, benchmark dependence, proxy metrics, and the gap between this thesis prototype and a
production CI system.
